% GENERATED, edit the French unit or the translation table
\mxunit{premiere-c06}{Differentiation, variation and optimisation}\label{C06}
\begin{keypoint}[Chapter path]
\textbf{Objectives.} Calculate a derivative on its domain; study variation and extrema; compare curves; optimise; use symmetries.

\textbf{Prerequisites.} Derivative at a point, algebraic calculation, sign of a quadratic, intervals.

\textbf{Yearly reference point.} 1--10 December: 8 hours, consolidation on the 15th and 17th\end{keypoint}
\section*{Diagnostic: six questions, 10 minutes}\label{C06-D}
Answer without consulting the lesson. Mark the questions for which help is needed. This diagnostic is not graded.
\begin{enumerate}
\item Derivative of $x^2$ at 3?
\item Domain of $\dfrac{1}{x-2}$?
\item Sign of $(x-1)(x+2)$?
\item Expand $(x^2+1)(x-2)$.
\item For $f(x)=x^2$, compare $f(-x)$ and $f(x)$.
\item If $f(1)=2$, $f\prime(1)=0$, what is the equation of the tangent?
\end{enumerate}
\paragraph{After the diagnostic.} Revisit the flagged topics with the teacher; redo a similar question, then explain the rule used.
\paragraph{Speaking and logic.} Present a variation table, distinguishing the sign of f from the sign of its derivative. Explain why the endpoints of an optimisation problem must be examined.
\mxnewpage
\section*{Investigate before formalising}\label{C06-RECH}
Can an increasing function have a horizontal tangent? Look for an example; can it have a strictly negative derivative at an interior point?
\paragraph{Organisation.} Individual investigation 5 min, comparison in pairs 5 min, class discussion 5 min. Keep a record of one attempt, even an unfinished one. A conjecture must then be justified.
\section{Derivative function and standard functions}\label{C06-U1}
\begin{keypoint}[Definition and results]
A function is differentiable on an interval when it is differentiable at each of its points; its derivative associates with x the number $f\prime(x)$. The standard derivatives are: constant $\mapsto0$; $x\mapsto1$; $x^2\mapsto2x$; $x^3\mapsto3x^2$; $\dfrac{1}{x}\mapsto -\dfrac{1}{x^2}$ for $x\ne0$; $\sqrt x\mapsto \dfrac{1}{2\sqrt x}$ for $x>0$. For $n\in\Z$, $(x^n)\prime=nx^{n-1}$ on its domain of differentiability; for $n=0$, the derivative is 0 directly.
\end{keypoint}
\paragraph{Status and justification.} Square: the rate is $2a+h$. Reciprocal: for $a\ne0$, $h\ne0$ and $a+h\ne0$, $\dfrac{[\dfrac{1}{a+h}-\dfrac{1}{a}]}{h}=-\dfrac{1}{[a(a+h)]}$, which tends to $-\dfrac{1}{a^2}$. The other formulas are accepted here.
\begin{example}For $f(x)=x^{-3}$, the domain is $\R\setminus\{0\}$ and $f\prime (x)=-3x^{-4}=-\dfrac{3}{x^4}$. For $g(x)=\sqrt x$, the function exists at 0, its standard derivative only for $x>0$.\end{example}\label{C06-EX1}
\begin{pitfall}The domain of differentiability may be smaller than the domain of definition.\end{pitfall}
\begin{application}Differentiate $x^5$, $x^{-2}$ and $\sqrt x$, stating the domains.\end{application}\label{C06-AP1}
\section{Sum, product, reciprocal and quotient}\label{C06-U2}
\begin{keypoint}[Definition and results]
For differentiable u and v, $(u+v)\prime=u\prime+v\prime$, $(\lambda u)\prime=\lambda u\prime$, $(uv)\prime=u\prime v+uv\prime$. Where $v\ne0$, $(\dfrac{1}{v})\prime =-\dfrac{v\prime}{v^2}$ and $(\dfrac{u}{v})\prime =\dfrac{u\prime v-uv\prime}{v^2}$.
\end{keypoint}
\paragraph{Status and justification.} For the product, add and subtract $u(a+h)v(a)$ in the numerator. The rate becomes $\dfrac{u(a+h)[v(a+h)-v(a)]}{h}+\dfrac{v(a)[u(a+h)-u(a)]}{h}$. Since a differentiable function is continuous (accepted), we obtain $u(a)v\prime(a)+v(a)u\prime(a)$. The other rules are accepted.
\begin{example}For $f(x)=\dfrac{x^2+1}{x-1}$, $x\ne1$, $f\prime (x)=\dfrac{[2x(x-1)-(x^2+1)]}{(x-1)^2}=\dfrac{x^2-2x-1}{(x-1)^2}$. Keeping the brackets avoids changing the sign of a single term instead of the whole numerator.\end{example}\label{C06-EX2}
\begin{pitfall}The derivative of a product is not the product of the derivatives; the derivative of a quotient is not their quotient.\end{pitfall}
\begin{application}Differentiate $(x^2+1)(3x-2)$, then $\dfrac{1}{x^2+1}$.\end{application}\label{C06-AP2}
\section{Sign of the derivative and extrema}\label{C06-U3}
\begin{keypoint}[Definition and results]
On an interval I, if $f\prime\geq0$, f is increasing; if $f\prime\leq0$, it is decreasing; if $f\prime=0$ everywhere, it is constant. A strictly positive derivative implies strict increase. If f has a local extremum at an interior point where it is differentiable, its derivative is zero there. The converse is false.
\end{keypoint}
\paragraph{Status and justification.} Variation theorems accepted. For an interior minimum at a, the rates on the right are positive or zero, those on the left negative or zero; if they converge to the same real number, it is 0. Same argument for a maximum.
\begin{example}$f(x)=x^3-3x$: $f\prime(x)=3(x-1)(x+1)$. f increases on $]-\infty;-1]$, decreases on $[-1;1]$, increases on $[1;+\infty[$; local maximum $f(-1)=2$, local minimum $f(1)=-2$. Neither of these extrema is global on $\R$.\end{example}\label{C06-EX3}
\mxinclude{figures/C06-U3}
\begin{pitfall}The theorem applies on an interval. A function can decrease on two separate intervals without decreasing on their union.\end{pitfall}
\begin{application}Study $f(x)=x^2-4x+1$ on $[0;3]$ and give its extrema on this interval.\end{application}\label{C06-AP3}
\section{Parity, inequalities and comparison}\label{C06-U4}
\begin{keypoint}[Definition and results]
On a domain symmetric about 0, f is even if $f(-x)=f(x)$ and odd if $f(-x)=-f(x)$. The curve of an even function is symmetric about the y-axis; that of an odd function about the origin. To compare f and g, we study the sign of $f-g$, possibly with its derivative.
\end{keypoint}
\paragraph{Status and justification.} The symmetries follow from the coordinates of the image points. For an inequality, a study of variation may give a minimum of the difference, then its sign; the derivative does not directly replace this difference.
\begin{example}For $x>0$, $h(x)=x+\dfrac{1}{x}-2$ has derivative $\dfrac{x^2-1}{x^2}$. It decreases up to 1 then increases; $h(1)=0$, so $x+\dfrac{1}{x}\geq 2$. Equality is obtained only at 1.\end{example}\label{C06-EX4}
\begin{pitfall}Check the domain before parity. The sign of $f\prime-g\prime$ tells us about the variation of $f-g$, not directly about its sign.\end{pitfall}
\begin{application}Study the parity of $x^4+x^2$ and $x^3-x$; prove $x^2+1\geq2x$.\end{application}\label{C06-AP4}
\section{Optimisation and Newton's method}\label{C06-U5}
\begin{keypoint}[Definition and results]
Optimising means defining the variable and its domain, expressing the quantity to optimise, then comparing all the candidate values, including the endpoints. Newton's method uses the tangent to approximate a solution of $f(x)=0$: $x_{n+1}=x_n-\dfrac{f(x_n)}{f}\prime (x_n)$ when $f\prime(x_n)\ne0$. It is used only in favourable cases.
\end{keypoint}
\paragraph{Status and justification.} Newton's formula comes from the intersection of $y=f(x_n)+f\prime(x_n)(x-x_n)$ with the horizontal axis. Convergence is not guaranteed in general; it is observed and checked in the case chosen.
\begin{example}For $f(x)=x^2-2$ with $x_0=1$, Newton gives $x_1=\dfrac{3}{2}$, $x_2=\dfrac{17}{12}$. The squares are 2.25 then about 2.00694: we approach $\sqrt2$. The program must check for a zero derivative and a maximum number of steps.\end{example}\label{C06-EX5}
\begin{pitfall}A numerical value produced by an algorithm is not a proof of a minimum or of convergence.\end{pitfall}
\begin{application}A rectangle of perimeter 12 m has width x. Determine the maximum area by differentiation.\end{application}\label{C06-AP5}
\section*{Exercise path}\label{C06-PARC}
Start with exercises 1 and 2. Then choose one of exercises 3 or 4 and do it without help. Exercises 5 to 7 are a reserve for transfer, investigation or programming: the teacher designates one. Exercise 8 is revisited after a delay.
\begingroup
\mxexercisehead{C06-E1}{A collection of derivatives}
Consider the functions\[f(x)=3x^4-2x+7,\quad g(x)=(x^2+2)(x-1),\quad h(x)=\dfrac{2x+1}{x+3}.\]
\begin{enumerate}[label=\arabic*.]
\item State the domain of definition of each of these functions.
\item Calculate their derivatives, naming the rule used each time. For $g$, first keep the two terms of the product rule; for $h$, write the numerator before simplifying it.
\item Check the result for $g$ by expanding its expression before differentiating.
\end{enumerate}
\mxexercisehead{C06-E2}{A complete table}
The function $f$ is defined on $[0;4]$ by $f(x)=x^3-6x^2+9x+1$. We look for its minimum and maximum values on this interval.
\begin{enumerate}[label=\arabic*.]
\item Calculate and factorise $f\prime(x)$, then study its sign on $[0;4]$.
\item Draw the variation table of $f$, giving the values at the endpoints and at the points where the derivative vanishes.
\item Give the minimum and the maximum of $f$ on the interval, as well as all the x-coordinates where they are reached.
\end{enumerate}
\mxexercisehead{C06-E3}{A quotient}
For every strictly positive real number $x$, let $f(x)=x+\dfrac4x$.
\begin{enumerate}[label=\arabic*.]
\item Study the variation of $f$ on $]0;+\infty[$. Present the factors that determine the sign of the derivative.
\item Deduce the smallest value of $x+\dfrac4x$, then write an inequality valid for all $x>0$. State exactly its equality case.
\end{enumerate}
\mxexercisehead{C06-E4}{Comparing two curves}
The functions $f(x)=x^3$ and $g(x)=x$ are defined on $\R$. Let $\mathcal C_f$ and $\mathcal C_g$ be their curves in the same frame.
\begin{enumerate}[label=\arabic*.]
\item Determine the points of intersection of these two curves.
\item Study their relative position on $\R$ using a suitable method. Express your result as the intervals where $\mathcal C_f$ is above or below $\mathcal C_g$.
\end{enumerate}
\mxexercisehead{C06-E5}{A box}
A box without a lid is made from a square sheet of side 12 cm. A square of side $x$ cm is cut from each of the four corners, then the remaining edges are folded up. The thickness of the sheet is neglected.
\begin{enumerate}[label=\arabic*.]
\item Draw the net and express the height and the two dimensions of the base of the box. State the values of $x$ that give a box of nonzero volume.
\item Express the volume $V(x)$ in cubic centimetres. Determine the dimensions that make this volume maximal.
\item Check the values at the two endpoints of the domain and give the exact maximum volume.
\end{enumerate}
\mxinclude{figures/C06-E5}
\mxexercisehead{C06-E6}{Stationary without an extremum}
A student claims: ``If the derivative vanishes at a point, the function necessarily has a local extremum at that point.'' We study the function $f(x)=x^3$ to test this claim.
\begin{enumerate}[label=\arabic*.]
\item Calculate $f\prime(0)$.
\item Compare with $f(0)$ the values of $f(x)$ for negative then positive real numbers $x$, as close to 0 as we like.
\item Conclude about the existence of a local extremum at 0 and about the claim proposed.
\end{enumerate}
\mxexercisehead{C06-E7}{Comparing strategies}
We wish to prove that, for every real number $x$, $x^2-6x+12\geq3$.
\begin{enumerate}[label=\arabic*.]
\item Give a proof by making a square appear.
\item Give a second proof by studying the variation of a function.
\item State the equality case. Compare the two methods: which requires the least calculation here?
\end{enumerate}
\mxexercisehead{C06-E8}{Delayed review}
The reciprocal function $f(x)=\dfrac1x$ has domain $D=\R\setminus\{0\}$. A student calculates $f\prime(x)=-\dfrac1{x^2}<0$ and concludes that $f$ is decreasing on all of $D$.
\begin{enumerate}[label=\arabic*.]
\item Identify the missing hypothesis in the application of the theorem linking the sign of the derivative and variation.
\item State the two intervals on which the conclusion is correct.
\item Using the numbers $-1$ and 1, construct a counterexample to the conclusion on all of $D$.
\end{enumerate}
\paragraph{Homework.} Do only the task designated at the end of the session, in the dated progression: 10--20 min at most. The exercises depend on the topics already studied, so some questions are deferred. Revisit exercise 8 one week after the chapter (10 min). If the allotted time is exceeded, note the difficulty and bring your attempts; do not pile up the whole reserve.
\endgroup
\mxneed{12}\section*{To do in the CQFD textbook}
These exercises are done in your textbook. The progression indicates the session; they replace the homework planned for that day (10 to 20 minutes in total).
\paragraph{CQFD, exercise 12 p.~156, questions 1 and 2.}\label{C06-M1} Calculate the difference, study its sign and translate this into the relative position of the curves.
\paragraph{CQFD, exercise 16 p.~156, questions f and g.}\label{C06-M2} For each function, give an interval on which it is both negative and increasing; justify the two properties separately.
\begingroup\small
\mxnewpage\section*{MCQ 1: understand and choose}\label{C06-Q1}
Five single-answer questions for practice. Choose one answer per question; then explain one distractor. This chapter format does not reproduce the official EAM format.
\begin{enumerate}
\item $(x^3)\prime$ equals
\begin{enumerate}[label=\Alph*.,itemsep=2pt]
\item $x^2$
\item $3x^2$
\item $3x$
\item $\dfrac{x^4}{4}$
\end{enumerate}
\item $(\dfrac{1}{x})\prime $ equals
\begin{enumerate}[label=\Alph*.,itemsep=2pt]
\item $-\dfrac{1}{x^2}$
\item $\dfrac{1}{x^2}$
\item $-\dfrac{1}{x}$
\item $1$
\end{enumerate}
\item $(uv)\prime$ equals
\begin{enumerate}[label=\Alph*.,itemsep=2pt]
\item $u\prime v\prime$
\item $u\prime+v\prime$
\item $uv\prime$
\item $u\prime v+uv\prime$
\end{enumerate}
\item $f\prime>0$ on an interval implies
\begin{enumerate}[label=\Alph*.,itemsep=2pt]
\item f positive
\item f negative
\item f strictly increasing
\item f constant
\end{enumerate}
\item $\dfrac{x^2+1}{x-1}$ is defined if
\begin{enumerate}[label=\Alph*.,itemsep=2pt]
\item $x>0$
\item $x\ne1$
\item $x\ne-1$
\item $x\ne0$
\end{enumerate}
\end{enumerate}
\endgroup
\begingroup\small
\mxneed{18}\section*{MCQ 2: understand and choose}\label{C06-Q2}
Five single-answer questions for practice. Choose one answer per question; then explain one distractor. This chapter format does not reproduce the official EAM format.
\begin{enumerate}
\item $x^3-x$ is
\begin{enumerate}[label=\Alph*.,itemsep=2pt]
\item even
\item constant
\item odd
\item neither even nor odd
\end{enumerate}
\item $f\prime(a)=0$ necessarily implies
\begin{enumerate}[label=\Alph*.,itemsep=2pt]
\item a horizontal tangent
\item a maximum
\item a minimum
\item f constant
\end{enumerate}
\item Derivative of $\sqrt x$:
\begin{enumerate}[label=\Alph*.,itemsep=2pt]
\item $2\sqrt x$
\item $\dfrac{1}{\sqrt x}$
\item $\dfrac{\sqrt x}{2}$
\item $\dfrac{1}{2\sqrt x}$, $x>0$
\end{enumerate}
\item Minimum of $x+\dfrac{1}{x}$ on $]0;+\infty[$:
\begin{enumerate}[label=\Alph*.,itemsep=2pt]
\item $0$
\item $2$
\item $1$
\item none
\end{enumerate}
\item Newton's method for $x^2-2$ is not defined with $x_0=$
\begin{enumerate}[label=\Alph*.,itemsep=2pt]
\item $1$
\item $2$
\item $0$
\item $-1$
\end{enumerate}
\end{enumerate}
\subsection*{True or false? Justify.}\label{C06-VF}
\begin{enumerate}
\item The derivative of a positive function is positive.
\item An interior extremum of a differentiable function gives a zero derivative.
\item The product of two increasing functions is increasing.
\end{enumerate}
\begin{keypoint}[Self-assessment]I differentiate on the right domain; I simplify without losing the forbidden values; I distinguish local from global extremum; I sometimes choose a factorisation rather than a derivative.

For each objective, indicate: achieved without help; achieved with help; to be reworked. Write one example and one mistake to avoid.\end{keypoint}
\endgroup
\ifmxteacher\mxinclude{corriges/C06}\fi
